% --- 题目 ---
\problem[P69-2 (25-3)]{微分方程 $\displaystyle xy' - y + x^2 \mathrm{e}^x = 0$ 满足 $\displaystyle y(1) = -\mathrm{e}$ 的解为 $\displaystyle y = \underline{\hspace{4em}}. $}
\ansat{281；【十年真题】 - 考点一：一阶微分方程的求解 - 2}

% --- 题目 ---
\problem[P69-3 (24-1,2)]{微分方程 $\displaystyle y' = \frac{1}{(x+y)^2}$ 满足条件 $\displaystyle y(1) = 0$ 的解为 \underline{\hspace{4em}}.}
\ansat{281；【十年真题】 - 考点一：一阶微分方程的求解 - 3}

% --- 题目 ---
\problem[P69-1 (17-2)]{微分方程 $\displaystyle y'' - 4y' + 8y = \mathrm{e}^{2x} (1 + \cos 2x)$ 的特解可设为 $\displaystyle y^* = (\quad)$
\begin{tasks}(2)
  \task $\displaystyle A \mathrm{e}^{2x} + \mathrm{e}^{2x} (B \cos 2x + C \sin 2x)$.
  \task $\displaystyle A x \mathrm{e}^{2x} + \mathrm{e}^{2x} (B \cos 2x + C \sin 2x)$.
  \task $\displaystyle A \mathrm{e}^{2x} + x \mathrm{e}^{2x} (B \cos 2x + C \sin 2x)$.
  \task $\displaystyle A x \mathrm{e}^{2x} + x \mathrm{e}^{2x} (B \cos 2x + C \sin 2x)$.
\end{tasks}}
\ansat{282；【十年真题】 - 考点二：高阶微分方程的求解 - 1}

% --- 题目 ---
\problem[P69-2 (22-2)]{微分方程 $\displaystyle y''' - 2y'' + 5y' = 0$ 的通解 $\displaystyle y(x) = \underline{\hspace{4em}}. $}
\ansat{282；【十年真题】 - 考点二：高阶微分方程的求解 - 2}

% --- 题目 ---
\problem[P69-4 (21-2)]{微分方程 $\displaystyle y''' - y = 0$ 的通解为 $\displaystyle y = \underline{\hspace{4em}}. $}
\ansat{282；【十年真题】 - 考点二：高阶微分方程的求解 - 4}

% --- 题目 ---
\problem[P69-5 (17-1)]{微分方程 $\displaystyle y'' + 2y' + 3y = 0$ 的通解为 $\displaystyle y = \underline{\hspace{4em}}. $}
\ansat{282；【十年真题】 - 考点二：高阶微分方程的求解 - 5}

% --- 题目 ---
\problem[P69-6 (24-2)]{设 $\displaystyle y(x)$ 为微分方程 $\displaystyle x^2 y'' + xy' - 9y = 0$ 满足条件 $\displaystyle \left.y\right|_{x=1}=2, \left.y'\right|_{x=1}=6$ 的解.
\begin{enumerate}
\item[(1)] 利用变换 $\displaystyle x=\mathrm{e}^t$ 将上述方程化为常系数线性方程, 并求 $\displaystyle y(x)$;
\item[(2)] 计算 $\displaystyle \int_1^2 y(x) \sqrt{4-x^2} \, \mathrm{d}x$.
\end{enumerate}}
\ansat{282；【十年真题】 - 考点二：高阶微分方程的求解 - 6}

% --- 题目 ---
\problem[P71-例(3) (90-2)]{微分方程 $\displaystyle x \ln x \, \mathrm{d}y + (y - \ln x) \, \mathrm{d}x = 0$ 满足条件 $\displaystyle \left.y\right|_{x=\mathrm{e}} = 1$ 的特解是 $\displaystyle y = \underline{\hspace{4em}}. $}
\ansat{71；【方法探究】 - 考点一：一阶微分方程的求解 - 例(3)}

% --- 题目 ---
\problem[P71-例1(2)]{微分方程 $\displaystyle y'' - y = \mathrm{e}^{-x} \sin x$ 的通解为 $\displaystyle y = \underline{\hspace{4em}}. $}
\ansat{71；【方法探究】 - 考点一：一阶微分方程的求解 - 例1(2)}

% --- 题目 ---
\problem[P74-例1 (10-2,3)]{设 $\displaystyle y_1, y_2$ 是一阶线性非齐次微分方程 $\displaystyle y'+p(x)y=q(x)$ 的两个特解, 若常数 $\displaystyle \lambda, \mu$ 使 $\displaystyle \lambda y_1 + \mu y_2$ 是该方程的解, $\displaystyle \lambda y_1 - \mu y_2$ 是该方程对应的齐次方程的解, 则 ( \quad )
\begin{tasks}(2)
  \task $\displaystyle \lambda = \frac{1}{2}, \mu = \frac{1}{2}$.
  \task $\displaystyle \lambda = -\frac{1}{2}, \mu = -\frac{1}{2}$.
  \task $\displaystyle \lambda = \frac{2}{3}, \mu = \frac{1}{3}$.
  \task $\displaystyle \lambda = \frac{2}{3}, \mu = \frac{2}{3}$.
\end{tasks}}
\ansat{74；【方法探究】 - 考点：已知微分方程的解的相关问题 - 例1}

% --- 题目 ---
\problem[P74-例2 (97-2)]{已知 $\displaystyle y_1=x\mathrm{e}^x + \mathrm{e}^{2x}$, $\displaystyle y_2=x\mathrm{e}^x + \mathrm{e}^{-x}$, $\displaystyle y_3=x\mathrm{e}^x + \mathrm{e}^{2x} - \mathrm{e}^{-x}$ 是某二阶非齐次线性微分方程的三个解, 则此微分方程为 \underline{\hspace{4em}}.}
\ansat{74；【方法探究】 - 考点：已知微分方程的解的相关问题 - 例2}

% --- 题目 ---
\problem[P75-4 (20-1)]{设函数 $\displaystyle f(x)$ 满足 $\displaystyle f''(x) + af'(x) + f(x) = 0$ ($\displaystyle a > 0$), 且 $\displaystyle f(0) = m, f'(0) = n$, 则 $\displaystyle \int_0^{+\infty} f(x) \mathrm{d}x = \underline{\hspace{4em}}. $}
\ansat{284；【十年真题】 - 考点：微分方程的应用 - 4}

% --- 题目 ---
\problem[P76-13 (20-3)]{设函数 $\displaystyle y=f(x)$ 满足
\[ y'' + 2y' + 5y = 0, \ f(0) = 1, \ f'(0) = -1. \]
\begin{enumerate}
\item[(1)] 求 $\displaystyle f(x)$ 的表达式;
\item[(2)] 设 $\displaystyle a_n = \int_{n\pi}^{+\infty} f(x) \mathrm{d}x$, 求 $\displaystyle \sum_{n=1}^{\infty} a_n$.
\end{enumerate}}
\ansat{286；【十年真题】 - 考点：微分方程的应用 - 13}

% --- 题目 ---
\problem[P77-16 (18-1)]{已知微分方程 $\displaystyle y'+y=f(x)$, 其中 $\displaystyle f(x)$ 是 $\displaystyle \mathbb{R}$ 上的连续函数.
\begin{enumerate}
\item[(1)] 若 $\displaystyle f(x)=x$, 求方程的通解;
\item[(2)] 若 $\displaystyle f(x)$ 是周期为 $\displaystyle T$ 的函数, 证明: 方程存在唯一的以 $\displaystyle T$ 为周期的解.
\end{enumerate}}
\begin{note}
    主要错的是 (2)
\end{note}
\ansat{286；【十年真题】 - 考点：微分方程的应用 - 16}

% --- 题目 ---
\problem[P77-17 (18-2)]{已知连续函数 $\displaystyle f(x)$ 满足
\[ \int_0^x f(t) \mathrm{d}t + \int_0^x t f(x-t) \mathrm{d}t = ax^2. \]
\begin{enumerate}
\item[(1)] 求 $\displaystyle f(x)$;
\item[(2)] 若 $\displaystyle f(x)$ 在区间 $\displaystyle [0,1]$ 上的平均值为 1, 求 $\displaystyle a$ 的值.
\end{enumerate}}
\begin{note}
    主要错的是 (1)
\end{note}
\ansat{286；【十年真题】 - 考点：微分方程的应用 - 17}

% --- 题目 ---
\problem[P77-19 (16-1)]{设函数 $\displaystyle y(x)$ 满足方程 $\displaystyle y'' + 2y' + ky = 0$, 其中 $\displaystyle 0 < k < 1$.
\begin{enumerate}
\item[(1)] 证明: 反常积分 $\displaystyle \int_0^{+\infty} y(x) \mathrm{d}x$ 收敛;
\item[(2)] 若 $\displaystyle y(0) = 1, y'(0) = 1$, 求 $\displaystyle \int_0^{+\infty} y(x) \mathrm{d}x$ 的值.
\end{enumerate}}
\ansat{286；【十年真题】 - 考点：微分方程的应用 - 19}
